\section{Introduction}
\label{sec:intro}

Real-time dynamic 3D reconstruction from video streams underpins a growing
list of applications --- robotic manipulation, autonomous driving, and
augmented and virtual reality --- that share a common requirement: the
reconstructed scene must carry \emph{absolute, metric scale}. A robot
grasping an object, a vehicle planning a lane change, and an AR system
anchoring a virtual asset all need to know where things are in meters, not
in a normalized canonical space. Yet the most successful recent feed-forward
dynamic 3D Gaussian Splatting (3DGS) methods~\cite{wu2026streamsplat}
operate on \emph{uncalibrated monocular} video and explicitly factor out
camera intrinsics and scale by working in an orthographic canonical space.
This makes them fast and broadly applicable, but leaves the user with a
scene whose metric scale is unspecified and must be recovered after the
fact, typically by matching against a separate SLAM trajectory or a known
object. The recovery step is fragile, introduces an extra source of error,
and defeats the purpose of a one-shot feed-forward model.

The physical source of metric scale is well understood: a calibrated stereo
rig with baseline $b$ and focal length $f$ recovers absolute depth via the
triangulation formula $Z = b f / d$, where $d$ is the disparity. Stereo
matching is a mature field with strong generalizing
backbones~\cite{xu2023igev}, and calibrated stereo cameras are ubiquitous
on modern robots and vehicles. This raises a natural question: \emph{can we
build a feed-forward dynamic 3DGS model that consumes a stereo video stream
and directly produces a metric-scale dynamic Gaussian scene, while
preserving the speed and online property of the monocular feed-forward
pipeline?}

We answer this affirmatively with \StereoSplat. Our framework reuses the
transformer encoder, probabilistic Gaussian decoder, and bidirectional
deformation field of \StreamSplat~\cite{wu2026streamsplat}, but makes three
targeted replacements. \textbf{First}, the monocular Depth-Anything depth
prior is replaced by a stereo-matching front end whose disparity output is
converted to metric depth through $Z = b f / d$. \textbf{Second}, the
fixed-camera orthographic rasterizer is replaced by a perspective
rasterizer that consumes the true intrinsics and per-frame poses, so that
the rendered images are geometrically consistent with the input cameras.
\textbf{Third}, we introduce a \emph{metric back-projection head} that
reinterprets the network's per-pixel position prediction as
sub-pixel/metric-depth offsets and lifts each Gaussian into the metric
world frame via the known intrinsics and pose. The dynamic decoder and
polynomial motion model of \StreamSplat are reused unchanged, except that
all quantities now live in meters.

Beyond reconstruction quality, dynamic 3DGS methods exhibit a family of
\emph{display} artifacts that limit their practical usefulness: temporal
flickering in static regions~\cite{yun2025compensating}, popping from
abrupt Gaussian appearance and disappearance, and high-frequency motion
blur under polynomial motion models~\cite{zhang2026fastgs}. We address
these with a streaming display module that wraps the rasterizer in a ring
buffer and applies two lightweight, training-free operators: a
motion-adaptive temporal exponential moving average (EMA) that strongly
smooths static regions while leaving high-motion regions intact, and an
opacity-stability filter that clamps per-frame opacity deltas to suppress
popping. The display module exposes a single \texttt{step()} API and
produces, alongside the rendered frame, diagnostic signals (motion energy,
popping score) that we use in our experiments to measure temporal
consistency.

Our contributions are:
\begin{itemize}[leftmargin=*,topsep=2pt,itemsep=1pt]
    \item We propose \StereoSplat, the first feed-forward framework that
    produces \emph{metric-scale} dynamic 3DGS from a stereo video stream,
    by anchoring scale to the stereo baseline and focal length.
    \item We introduce a metric back-projection head and a perspective
    rasterization path that integrate seamlessly with the existing
    \StreamSplat architecture, reusing its probabilistic and dynamic
    components.
    \item We identify and address the display artifacts of dynamic GS with
    a streaming display module combining motion-adaptive EMA and
    opacity-stability filtering.
    \item We conduct a comprehensive experimental study on TartanAir,
    DSEC, and KITTI Stereo, including a five-axis ablation, baseline
    comparisons, and qualitative visualizations, demonstrating both metric
    scale recovery and competitive photometric quality.
\end{itemize}
